\chapter[Race and Its Real Consequences]{Race: Real Consequences, No Biological Hierarchy}
\section{A Pencil's Verdict}
First, pick up an object. It is neither an antiquity nor a weapon. It may be lying in your pencil case right now: an ordinary wooden pencil.

In mid-twentieth-century South Africa, a pencil like this could determine a person's entire life.

The procedure sounds like a vicious joke today. An official would insert a pencil into the hair of the person being examined and ask them to lower or shake their head. If the pencil fell out, the hair was straight and soft enough: the person might be registered as "Coloured," or, with greater luck, even "White." If it stayed lodged in tightly curled hair, the person was "Black." This is no folk tale. It was one of the tests actually used under apartheid, and it had a name: the "pencil test." Whether a pencil fell determined which neighborhood you could live in, whom you could marry, which school your children could attend, and which cemetery could receive you after death.

Weigh this pencil in your hand. It measures neither intelligence nor character, not even degrees of ancestry: only the curliness of hair. Yet its consequences could outweigh any written judgment. More absurd still, a pencil might fall during one test and remain in place during the next, potentially rewriting the same person's identity. Siblings in one family really were assigned to different "races." Legally, they now belonged to different "racial types" and could no longer live in the same neighborhood or attend the same school.

The question this chapter asks is hidden inside that pencil. If the boundaries of "race" are so absurd that a pencil must draw them, how can they possess the power to rewrite lives? People often say that "race is not a biological hierarchy," yet which of its products---segregation, poverty, violence---is unreal? How can something that "does not exist in nature" have consequences so real?

\section{An Office in Pretoria}
In 1950, South Africa passed the Population Registration Act. In essence, everyone in the country had to be registered in a racial category, with identity documents stating in black and white whether they were "White," "Coloured," or "Black." Once the law existed, somebody had to enforce it. Hence the Race Classification Board.

The following scene is a composite drawn from many real cases.

The setting: a small room in a government office building in Pretoria. It smells of smoke and old paper. Three officials sit behind a desk covered with forms and a magnifying glass---yes, a magnifying glass, for inspecting hair texture and facial features in identity photographs. A family stands before the desk: father, mother, two children. They wear their best clothes, because today's "interview" will decide the entire family's identity.

Their registered category is in question. The family was originally classified as "Coloured"---a sprawling South African category between "White" and "Black"---but a neighbor has reported that they "look more like Blacks." This is not scientific curiosity. If the family is reclassified as "Black," they must leave their present neighborhood, legally reserved for another race, and their house will become vacant.

What the officials do next captures the essence of this farce. They inspect the family's skin color---but skin changes with the seasons; does a suntan count? They examine hair, and the pencil makes its entrance. They study the widths of facial features, even asking family members to turn, walk, and speak, listening for a "Black accent." They discuss whom the family socializes with and which church they attend, because "social reputation" is also a criterion. If you live "like a White person," you can apply to be "reclassified" as White.

Notice a crucial detail: throughout this process, not a drop of blood is drawn and not a page of genealogical records consulted. Even if officials wanted such records, they would have nowhere to start, since the ancestors of most people in this country were never accurately registered. The "race" the board judges is not discovered inside the body. It is assembled from eyes, hair, accents, and neighbors' remarks.

Such classification hearings continued for more than forty years, from 1950 until the law's repeal in 1991. Thousands applied to change categories. Some successfully "became White," and their lives switched tracks. Others were "made Black," losing homes and jobs and being forced with their families into barren so-called "homelands." One tick on a form sent destinies down different paths.

We have visited the scene. Carry the smell of that room with you as we go on.

\section{Visible Differences, Inexplicable Boundaries}
Now let us lay out the conflict without answering it yet.

On one side stands an undeniable fact: people really do look different. Skin can be darker or lighter, hair straight or curly, nose bridges higher or lower. These differences are written on bodies, visible to the naked eye, and clearly associated with geographical ancestry. People whose ancestors lived for generations in strongly sunlit equatorial regions generally have darker skin; those whose ancestors lived at high latitudes with weaker sunlight generally have lighter skin. To claim that "there are no biological differences between people" is to deny the obvious. No serious geneticist would say it.

On the other side stands an equally undeniable fact: the "racial" boundaries drawn throughout human history are absurdly fragile. A pencil can shift them. Children of the same parents can belong to two "races." The United States long followed the so-called "one-drop rule": even one-eighth or one-sixteenth African ancestry could suffice for registration as Black. In Brazil, the same person might fall under any of more than a dozen skin-color labels. When Brazil's 1976 census let people describe their own color, it received over a hundred answers, ranging from "toast-colored" to "coffee with milk." The same body could be "Coloured" in South Africa, "Black" in the United States, and "light-colored" in Brazil. If the category really resided inside the body, why would it change at a national border?

Here is the real conflict: \textbf{the differences are real, but the knife used to divide them is plainly human-made.} So---

\textbf{What exactly is "race"?} Is it a natural hierarchy inside bodies, or a system of classification invented by society? If the latter, why do so many people around the world firmly believe the former? If it was invented, why are its consequences---segregation, enslavement, massacres, disparities of wealth---real enough to measure?

Do not choose an answer yet. Consider a simpler example. If your classmates line up by height, the line is real. But are "tall people" and "short people" two species that already exist in nature? Clearly not. Height differences are real; the division into "two kinds of people" is human-made. That distinction is this chapter's key. Be careful, though: race is far more complicated than height, because a whole history and institutional system rides on that knife, and has done so for centuries. Let us continue.

\section{Who Says Race Is Natural, and Who Says No?}
The idea that "race is a biological hierarchy" did not spring from nowhere. Generations of powerful people built it, brick by brick, for specific interests. Listen to what they said, and to those pressed beneath them.

\textbf{Slaveholders and their defenders.} In the eighteenth and nineteenth centuries, the Atlantic slave trade transported more than ten million Africans to the Americas. Notice the sequence: first came an extraordinarily profitable business; only then was a theory needed to defend it. Early defenses were mainly religious. Africans were called "descendants of Ham," supposedly destined for servitude---an appropriation of a biblical story. American independence brought a problem: how could a country proclaiming that "all men are created equal" retain slavery? The defense was upgraded and dressed in scientific clothing. Doctors declared African lung capacity, pain sensitivity, and brains "different from" those of Whites. Others claimed Black people were "unsuited" to freedom, making slavery virtually an act of kindness. That was the gist. The central maneuver was to rewrite an economic institution as a law of nature: I do not want to enslave you; nature has made you fit to be enslaved.

\textbf{Scientists did not speak with one voice.} We should be honest: nineteenth-century science was not a monolithic chorus of villains. Some researchers sincerely believed they were conducting objective inquiry. Carl Linnaeus of Sweden---the man who established binomial nomenclature for plants and animals---also divided humans into several "varieties" in his Systema Naturae, attaching assessments of temperament and morality: one variety was "irritable" and "stubborn," another "gentle" and "governed by laws." Today this looks like prejudice written into taxonomy. Then it counted as respectable European scholarship. In the late eighteenth century, Germany's Johann Friedrich Blumenbach proposed five categories and coined "Caucasian." Intriguingly, Blumenbach himself opposed ranking humanity and collected works by authors of African descent to demonstrate their talents. But once his five drawers circulated, others filled them with "ranks." However restrained a scientific tool and its maker might be, they could not restrain a society eager for the conclusion it wanted.

\textbf{The people being classified.} Frederick Douglass was born enslaved, taught himself to read, escaped north, and became one of nineteenth-century America's most powerful abolitionist speakers. Slaveholders' theories called Africans "naturally dull and lacking reason." Douglass's very presence at a lectern was a living refutation. More radically, he exposed how the claim of "natural destiny" was produced. In his autobiography, he detailed how an infant---himself---was systematically "manufactured" into a slave: separation from his mother, prohibition of literacy, a will broken by whipping. Slaves were made, not born. We will see the full weight of that statement when we read the evidence.

\textbf{Give the "seeing is believing" argument its strongest case.} Try an exercise: state as fully as possible a position that still exists on the opposing side today. It says: "Forget history; use your eyes. East Asians, Europeans, and West Africans really look different. An ancestry-testing company can take a tube of saliva and tell you what percentage of East Asian ancestry you have. Forensic specialists can infer a dead person's likely ethnicity from a skeleton. If race is not biological fact, how could these techniques work?"

This is not a weak argument, and it deserves a serious response. Ancestry testing really works: certain stretches of your genes do carry geographical information. In the tens of thousands of years since humans left Africa, marriage between populations has never been entirely unrestricted, and geographical separation has left statistical traces in genes. Forensic specialists really can infer probable ancestry from bones. There is no shame in acknowledging this.

But notice the argument's quiet change of concepts halfway through. It demonstrates that "\textbf{ancestry leaves biological traces}," then asks you to accept that "\textbf{race is a biological hierarchy}." The distance between those statements is the entire distance this chapter travels. Ancestral traces form a continuous, statistical, gradually changing pattern---like temperature changing gradually from south to north, with no latitude engraved "tropical race begins here." "Race," by contrast, claims there are several discrete boxes arranged in ranks. An ancestry company's report gives "percentages," never "your human rank," because the former is in the genes and the latter is not. Genes contain gradients, not boxes and ranks. The pencil test exposes exactly this: if the boxes existed naturally, why would a pencil be needed to decide who belongs in which?

\section{A Bridge for Thought: Three Layers of Race}
Here is a tool to take away. Whenever you encounter an argument about "race"---in the news, a textbook, or dinner-table conversation---first separate three layers that have been compressed together.

\textbf{Layer one: biological differences in ancestry.} People's genes do differ, and the frequencies of some differences correlate with geographical ancestry. This is genetics' territory. But examine its actual shape. In the 1970s, geneticist Richard Lewontin made a classic calculation, distributing human genetic variation across individuals, populations, and continents. His conclusion surprised many: about 85 percent of all human genetic variation occurred \textbf{between individuals within the same population}; only a small remainder lay between major "races." Later large-scale sequencing repeatedly confirmed this order of magnitude. In other words, choose any two people from the same village and their genetic differences greatly exceed the difference between "the average European" and "the average African." Another frequently cited fact is that any two human genomes are approximately 99.9\% identical. Biological differences exist, but their distribution simply does not fit the mold of "a few clearly ranked races."

\textbf{Layer two: social systems of classification.} This layer is human-made: how many boxes, what names, who goes where, and whether the boxes are ranked. It varies across places and periods. America's "one-drop rule" counted someone with one-eighth African ancestry as Black, serving to ensure that slaves' children were born slaves: classification followed property rights. Brazil's hundred-plus color terms reflected the continuum created by centuries of intermarriage; with too many shades to box, classification lost some of its sharpness as an instrument of ranking. South Africa's pencil test manufactured boundaries for an apartheid state. Three countries cut the same biological raw material into three entirely different sets of "races." A simple clue identifies classifications at this layer: change the country or the period, and they change.

\textbf{Layer three: institutional consequences.} Once classifications enter laws, schools, banks, and hospitals, they produce real consequences: who can buy a home, which schools children attend, whose average lifespan is several years longer. This layer is real too, but in an entirely different way from layer one: its reality is social. Money is socially real in the same sense. The paper itself has little value, but collective recognition lets it buy houses and land.

Most confusion comes from collapsing these three layers into one. Slaveholders disguised layer two as layer one: "Nature gave us our classification." Some naive rebuttals deny layer one: "Since classification is human-made, people have no biological differences." That is equally wrong, and ancestry testing immediately disproves it. The right approach lets each layer answer for itself: acknowledge biological differences; reject natural ranks; recognize human-made classifications; never dismiss their consequences merely because they are "human-made."

Next time you hear people arguing over "whether race exists," first ask: which layer do you mean? Eight out of ten arguments may stop immediately, because the speakers were not discussing the same layer at all.

\section{Two Statements by the Same Person}
Now read primary materials. Begin with a passage known around the world, from the American Declaration of Independence of 1776:

\begin{quote}
"We hold these truths to be self-evident: all people are born equal, and their Creator has endowed them with certain rights that cannot be taken away, including life, liberty, and the pursuit of happiness."
\end{quote}
Thomas Jefferson was the Declaration's principal drafter. Now turn to another book by the same person. Discussing people of African descent in Notes on the State of Virginia (1785), Jefferson expressed a "suspicion": in essence, he suspected that Blacks were inferior to Whites in reasoning and imagination, presenting this as a "difference made by nature." He acknowledged that it was only a suspicion and the evidence inadequate. But he wrote it down, and the book circulated widely.

Put these materials together and the central fracture of the "racial system" appears. One man could write "all people are born equal" in one document, rank people in another, and enslave hundreds over his lifetime. "Hypocrisy" alone does not explain it. A way of thinking was at work: \textbf{quietly shrinking the scope of "everyone."} Once slaves fell outside "everyone," the proclamation of equality remained undisturbed. Racial classification reveals its function here: a fence keeping certain people outside the "everyone" in "everyone is born equal."

Now consider testimony from someone inside that fence. In Narrative of the Life of Frederick Douglass (1845), Douglass described finally resisting the overseer whose specialty was "breaking" slaves when he was sixteen. They fought for two hours, and afterward the overseer never dared touch him again. He wrote (rendered here from the Chinese translation of the English original):

\begin{quote}
"You have seen how a person was made into a slave; now you will see how a slave was made back into a person."
\end{quote}
Read the sentence's structure closely. Douglass does not say, "I recovered my natural freedom." He speaks of two acts of making: a person \textbf{is made} a slave; a slave \textbf{is made} a person. He demotes slavery from a "natural condition" to a \textbf{manufacturing process}, with stages, tools---separating mother and child, banning literacy, whipping---and a purpose. The opposite of this "manufacture" is not "nature" but another kind of making: resistance, literacy, narration. The autobiography caused a sensation in Europe and America. Many readers could not believe a formerly enslaved person could write such prose. Their astonishment itself revealed how thoroughly the theory of "natural dullness" had conditioned them.

One document from the summit of power, another from inside the fence: together they say the same thing. Racial hierarchy was not "discovered" but \textbf{written}---into laws, books, and forms---before it grew into people's vision, making observers think it had always been there.

\section{How Was This Classification Made?}
The factual picture is clear: the modern racial system took shape with the Atlantic slave trade and European colonial expansion, acquired a pseudoscientific gilding in the nineteenth century, and became institutionalized in national laws in the twentieth. These points are supported by evidence and are not seriously disputed. What remains contested is explanation: why did humans build this system, and work so hard at it? Here are three explanations. They are not equivalent; each has leverage and blind spots of its own.

\textbf{Explanation one: interests---enslavement first, theory afterward.}

This account sees racism fundamentally as a machine for justification. Its clue is chronological: Europeans traded Africans on a large scale before systematic racial theories appeared. Each upgrade in justification precisely matched a new demand: religion when religious defenses worked, "science" stepping in after declarations of equality. Many economic historians follow this line. Prejudice did not create slavery; slavery created prejudice. Classification followed profit. Another piece of evidence: early Irish immigrants to North America were also "racialized," portrayed in newspapers as apes and described as inferior. As they established themselves economically and politically, they gradually "became White." If inferiority were innate, how could money and votes cure it? This sharp explanation has substantial evidence. Its limitation is its inability to explain why prejudice outlives the interests behind it. More than a century after slavery's abolition, its classifications still affect hiring and lending. The justification machine's owner is dead; why does the machine keep running?

\textbf{Explanation two: cognition---human brains love categories.}

This psychological account begins with the brain as a classification machine, sorting everything, including people, into compartments. People also naturally favor their own compartment, a tendency psychologists call "in-group favoritism." Even when experiments randomly divide people into "red" and "blue" teams by lot, participants begin favoring teammates within minutes. On this account, racial classification is a malignant instance of a general human instinct. Interests did not create prejudice; prejudice came in the brain's factory settings, and interests merely flipped the switch. This explains the embers that the interests account cannot: the machine keeps running after its owner dies because cognition supplies its fuel. Its limits are equally clear. If classification is purely instinctive, why do cultures classify so differently? Instinct explains why people categorize, not why they use skin color, arrange categories in ranks, or become so brutally committed to them in particular centuries. Attributing racism to instinct also risks sliding into "Nothing can be done; that is human nature"---exactly the conclusion perpetrators most want to hear.

\textbf{Explanation three: authority---pseudoscience gave prejudice a diploma.}

This account focuses on nineteenth-century science. Racial theory was more poisonous than older prejudices because it bore the highest seal of intellectual authority then available. In the 1830s and 1840s, American physician Samuel George Morton collected more than a thousand human skulls, used craniometry to measure brain capacity, and ranked "Caucasians largest, Africans smallest." Books including Crania Americana (1839) were treated as rigorous empirical science. Later reviews found that errors in his data handling almost always favored his expected conclusions: choices of samples and discarded data helped the desired result along. This explains why racial systems could stride confidently into schools, museums, and government: people in white coats had vouched for them. Its weakness is that it makes scientists sound too much like masterminds. More often, they too inhabited their era's prejudices, largely "verifying" their society's common sense rather than inventing malice from nothing. Authority explains racial theory's credibility, but not why that credibility was so eagerly misused.

Each explanation grasps one stretch of a causal chain. Interests supplied the motive, cognition the fuel, and authority the diploma. In actual history, all three chains intertwined. That also explains the system's persistence: remove any one, and the other two can keep it running.

\section[The Unreal Has Real Consequences]{A Deeper Challenge: How Can the Unreal Have Real Consequences?}
We now reach the chapter's hardest theoretical question. We have called racial classifications "social constructions": they are not things waiting to be discovered in nature, but products of human society. Yet "constructed" does not mean "fictional," much less "unimportant." What is a "constructed reality"? Begin with an everyday example: money.

A sheet of paper printed with a portrait and numbers can buy rice, houses, and human labor. Where does its value reside? Not inside the paper: a counterfeit detector checks security features, not value. Value resides in shared recognition. You accept it because you believe others will; because everyone believes this, it actually works. Money is a classic \textbf{social fact}: constituted entirely by belief, yet with consequences harder than most natural facts. Try withholding rent and telling your landlord that "money is only a social construction."

Sociologists W. I. Thomas and Dorothy Thomas wrote a much-quoted sentence in The Child in America (1928), rendered here from the Chinese translation of the English original:

\begin{quote}
"If people define situations as real, their consequences will be real."
\end{quote}
It could almost have been written for this chapter. Racial classifications are "defined as real" through laws, forms, pencils, and skull measurements. Their consequences are undeniably real: people assigned to inferior boxes cannot enter good schools, obtain loans, or live in desirable neighborhoods, and have shorter life expectancies. The definition is false; the consequences are real. Here is the solution to the "racial paradox": \textbf{race does not exist as a biological hierarchy, but does exist as a social machine; that machine's components happen to be people, who bear the weight of classification throughout their lives.}

Consider another analogy geneticists favor, one that firmly separates "real consequences" from "natural causes." Randomly divide a batch of maize seeds harvested from one field into two groups. Plant one in fertile soil and the other in poor soil. The fertile-field plants grow considerably taller on average. What causes differences between tall and short plants within the fertile field? Mainly genetic differences: water, nutrients, and sunlight are the same, so seeds determine height. What causes the average difference between fertile and poor fields? Almost entirely the environment---the soil. Here is the crucial point: \textbf{a genetic explanation for differences within each group never establishes a genetic explanation for differences between groups.} The seed groups share a gene pool, yet their average heights differ dramatically. Replace "fertile and poor fields" with "groups enjoying quality education and health care" and "groups institutionally deprived of them." You have the methodological baseline for understanding statistical differences between groups. When average achievement, income, or lifespan differs, the first question is not "How different are their genes?" but "How different is their soil?"

With this theory, "institutional discrimination produces measurable consequences" becomes a statement you can measure concretely. Consider two repeatedly documented mechanisms. First, American "redlining" in the 1930s: federally supported agencies marked predominantly Black neighborhoods in red as "unsuitable for lending," and banks refused loans accordingly. Decades of property-wealth accumulation were cut off, leaving effects still measurable in house prices and school-district disparities more than half a century later. Second, from 1932 to 1972, American public-health institutions concealed diagnoses and withheld treatment from hundreds of Black men with syphilis at Tuskegee simply to "observe the disease's course." Forty years: institutional discrimination measured directly in years and lives. These examples are not invitations to stare at suffering. Their point is mechanism: how consequences travel, link by link, from "definition" into bodies. Understanding mechanisms does not excuse perpetrators. On the contrary, only by seeing them can we identify where responsibility lies.

An easily misunderstood counterexample belongs here, lest you misread the chapter. Hearing that "race is not a biological hierarchy," some immediately conclude that every medical distinction by race is pseudoscience. Not so: reality is finer-grained. Some genetic variants really do differ in frequency by ancestry. Variants associated with sickle-cell anemia occur more frequently where malaria has long been endemic, as an adaptation against malaria. Those regions include West Africa, but also the Mediterranean, the Middle East, and India. The box labeled "Black disease" therefore fails to match reality: Greeks and residents of parts of Italy have substantial carrier rates, while many southern African populations classified as "Black" have low rates. What medicine should actually use is \textbf{ancestry and specific risk factors}, not racial boxes---the direction of recent medical changes. The lesson cuts both ways: do not mistake racial boxes for biology, and do not use "race is constructed" to erase real ancestral variation. The three-layer distinction has helped us again.

\section[Algorithms, Clinics, and Identity Cards]{Today's Echoes: Algorithms, Clinics, and Identity Cards}
This old lesson is taught again every day, in new classrooms.

\textbf{At school and online.} You have probably encountered this argument: "Why discuss race or ethnicity in this day and age? Just ignore skin color and treat everyone equally." This "color-blind" position sounds thoroughly decent. Apply our tool: it denies layer two, classification, hoping thereby to cancel layer three, consequences. But consequences do not disappear when nobody mentions them. Redlined neighborhoods do not catch up in property values because race goes unspoken; applications filtered out by prejudice do not receive callbacks because interviewers call themselves "color-blind." Color-blind goodwill often preserves precisely the inequality it wishes to eliminate, because it discards the measuring instrument that reveals the problem. This does not mean fastening people permanently into categories. It means acknowledging that a bad classification's accounts remain outstanding before dismantling it.

\textbf{In technology.} Algorithms and artificial intelligence increasingly perform work once done by officials and pencils: screening applications, approving loans, assessing "reoffending risk." An algorithm does not test hair with a pencil, but it consumes historical data soaked in the consequences of old classifications. A screening model trained on decades of admission records learns past discrimination as a "pattern," then executes it dispassionately at scale. Prejudice has been upgraded from pencil to code. Both claim objectivity.

\textbf{In medicine.} Consider a concrete change now under way. Common formulas estimating kidney function, or eGFR, long included a "race coefficient": a correction factor for "Black" patients, assuming greater muscle mass and higher blood creatinine. After reconsideration, kidney specialists in the United States and several other countries judged this crude racial box more harmful than helpful and, from 2021, promoted race-free equations. Every detail of this debate falls within our tool's reach. Nobody denies statistical differences in average creatinine levels between populations, layer one. The question is whether approximating ancestry and muscle mass with racial boxes harms individual patients---the crudeness of layer two and the consequences of layer three. It can, for example, overestimate some patients' kidney function and delay their place on transplant waiting lists. Science here is not denying differences; it is \textbf{replacing crude boxes with more accurate tools}.

\textbf{On the world map.} Racial systems have never belonged to one country alone. Rwanda's distinction between Hutu and Tutsi was originally closer to a fluid social status. In the 1930s, Belgian colonizers turned it into fixed categories printed on identity cards, distributing education and official posts accordingly. Half a century later, those cards became passes to death during the 1994 genocide: in roughly one hundred days, hundreds of thousands were killed, commonly estimated at about eight hundred thousand, overwhelmingly Tutsi. Remember that the dead are not footnotes to "ethnic hatred": each had a name and a life. Most killers, too, were ordinary people trained by the classification machine. Understanding its assembly is precisely how we prevent another from leaving the factory. India offers another example. Caste is not "race," but it is a similar device: a supposedly divinely ordained hierarchy joined to separation in marriage, occupation, and residence. More than two thousand years of operation have likewise left measurable inequalities today. Together, these examples reveal our real question: not "How does one race regard another?" but \textbf{why do human societies repeatedly manufacture classifications of "natural hierarchy," and why do people pay for them with entire lives?}

\section[Should the Form Keep That Field?]{Your Question: Should the Form Keep That Field?}
To finish, the pen is yours.

Many countries still ask for "race" or "ethnicity" on census, school, and employment forms. This chapter offers two opposing positions, each with genuine reasons.

Those who would remove the field argue: if classifications are constructed and harmful, and every box tells a lie, stop lending state machinery to keep them alive. Every completed form reclassifies; every tick reinforces a box. Only by ceasing to count can we eventually cease to distinguish.

Those who would retain it argue: consequences do not disappear because we stop looking. Without that field, there are no data. Without data, disparities in school funding, disease rates, and lifespans become invisible---and invisible inequality cannot be corrected. Race is constructed, but precisely for that reason we need statistics to keep watch over its construction site.

Now answer: \textbf{should forms retain the "race/ethnicity" field?} Why?

Before answering, take one final inventory. You know biological differences exist but natural ranks do not: the three-layer distinction. You know classifications are human-made and change with interests and eras: pencils, one-drop rules, a hundred color names. You know how false definitions grow real consequences: the Thomas theorem, fertile and poor fields. You have also seen the costs of both "not looking" and "keeping watch."

There is no standard answer. But notice something gentle: the process of finding your reasons is itself the most complete rejection of that pencil. The pencil needs one second to pass judgment on a person. You, a reader unwilling to let any box do your thinking, have devoted a whole chapter.
